% ============================================================
%  Georg Brandes, Samlede Skrifter. Trettende Bind (Supplementbind).
%  Kjøbenhavn: Gyldendalske Boghandels Forlag (F. Hegel & Søn), 1903.
%  Five pieces on Rasmus Nielsen and his school, 1867.
%  TRANSLATION (English), printed pp. 85--109.
%  Translated from transcription.tex (Danish) in this directory, whose
%  header records how its text was established.  Method:
%  ../../../../TRANSLATION-PLAYBOOK.md.  The volume is Brandes's 1903
%  revision of pieces first printed 1865--69; the translation follows the
%  1903 text.
%
%  ---- REGISTER ----
%  Readable modern English.  Brandes's long periods are broken where English
%  will not carry them, but nothing is softened: the polemic is sarcastic,
%  and where the Danish is ironic the English is ironic.
%
%  ---- TERMINOLOGY ----
%  Aligned with texts/brandes/dualismen/translation.tex and the Nielsen and
%  Brøchner translations in this repo, since the pieces belong to the same
%  controversy:
%    Tro / Viden          = faith / knowledge   (Tro og Viden = faith and
%                           knowledge)
%    Erkendelse           = cognition (NOT "knowledge", reserved for Viden)
%    Videnskab            = science (Wissenschaft; never "natural science",
%                           for which Brandes writes Naturvidenskab)
%    Fagvidenskaberne     = the special sciences
%    absolut Uensartethed = absolute heterogeneity
%    Aand / aandelig      = spirit / spiritual;  Fornuft = reason;
%    Forstand             = understanding;  Begreb = concept
%    Samvittighed         = conscience;  Selvansvarlighed = self-responsibility
%    Livsanskuelse        = view of life;  Verdensanskuelse = world view
%    Dannelse             = cultivation;  Folkeaand = national spirit
%    Aabenbaring          = revelation;  Under = miracle
%  Danish work-titles stay in Danish, as Brandes cites them (Grundideernes
%  Logik, Om den gode Vilje, Fædrelandet, Dagbladet, ...).
%
%  ---- CONVENTIONS ----
%  - Quotation marks: the volume's single guillemets ‹...› -> ``...''.
%  - Emphasis: italic is the volume's only emphasis device; every \emph{}
%    span of the transcription is carried over as \emph{}, including the
%    Latin and other foreign phrases, which stay untranslated.
%  - Footnotes \fnmark{*)}{...} at the same anchor.
%  - The ɔ: id-est mark -> "i.e."
%  - Page markers \opage{N} and "% --- p. N ---" at the same joints as in
%    transcription.tex, so the two can be read side by side.
%  - Heads: a book under review keeps its Danish title, with an English
%    gloss in brackets; Brandes's own titles are translated.  The table of
%    contents gives the English of the volume's Indhold titles.
%
%  ---- HOW THE TRANSLATION WAS CHECKED (2026-09-22) ----
%  (1) Mechanical audit (tr_audit.py): per printed page, the same \opage
%  joints, \emph spans, footnotes, heads, rules and paragraph breaks as
%  transcription.tex, and an English/Danish word ratio within 0.95--1.60:
%  25 pages, 0 flagged.  (2) Independent review: a reader who had not
%  translated these pages read every sentence of the Danish against the
%  English.  No omission found; 8 corrections at 11 sites (e.g. `Troen paa
%  Dæmoner' "belief", not "faith", p. 87; `aandeligt stavnsbundne'
%  "intellectually bound to the soil", pp. 96, 104; `Aand' "mind" of a
%  thinker's powers, pp. 107--108; `synes at maatte' hedge restored, p. 108).
%  One sentence is left open: p. 100 `at naa Højdepunktet i Rationalismens
%  Betragtning af det Kristelige' is rendered literally, keeping its
%  ambiguity; the comment at the site gives both readings.
%
%  ---- DEFECTS: DANISH AS PRINTED, ENGLISH CORRECT ----
%  The transcription is diplomatic and carries every printer's error; the
%  translation renders the SENSE and logs the divergence in a % comment at
%  the site.  The English quotation marks balance throughout.
% ============================================================

\documentclass[12pt,a4paper,openany]{book}
\usepackage[utf8]{inputenc}
\usepackage[T1]{fontenc}
\usepackage{libertinus}
\usepackage{libertinust1math}
\usepackage{textalpha}
\usepackage{graphicx}
\usepackage[english]{babel}
\usepackage[protrusion=true,expansion=true]{microtype}
\usepackage{setspace}
\usepackage[margin=1.2in]{geometry}
\usepackage{enumitem}
\usepackage{fancyhdr}
\setlength{\headheight}{28pt}

% Footnotes as the 1903 printing sets them: *), set at each site with
% \fnmark{*)}{...}, exactly as in transcription.tex.
\renewcommand{\thefootnote}{*)}
\newcommand{\fnmark}[2]{\renewcommand{\thefootnote}{#1}\footnote{#2}%
  \renewcommand{\thefootnote}{*)}}

% A piece's head, mirroring transcription.tex.  #1 = the English of the
% title as the volume's Indhold gives it (table of contents, running head);
% #2 = the head as printed, translated (a book under review keeps its Danish
% title, glossed in brackets).  The short rule under every head is printed.
\newcommand{\piecehead}[2]{\clearpage\phantomsection
  \addcontentsline{toc}{section}{#1}\markboth{#1}{#1}%
  \begin{center}\large #2\\[6pt]\rule{2cm}{0.4pt}\end{center}\par\medskip}
\newcommand{\piecerule}{\par\begin{center}\rule{1.5cm}{0.4pt}\end{center}\par}

\usepackage{hyperref}
\hypersetup{pdftitle={Brandes, Collected Writings XIII: Reviews and Replies on Rasmus Nielsen (1867)}, pdfauthor={Georg Brandes}, hidelinks}

\pagestyle{fancy}
\fancyhf{}
\fancyhead[LE,RO]{\thepage}
\fancyhead[RE]{\textit{Collected Writings XIII (1903)}}
\fancyhead[LO]{\textit{\leftmark}}
\setstretch{1.3}

\title{\textbf{Reviews and Replies on Rasmus Nielsen (1867)}\\[8pt]\large Georg Brandes, \textit{Samlede Skrifter} [Collected Writings], vol.~XIII (Supplementbind),\\ pp.~85--109}
\author{}
\date{Copenhagen: Gyldendal (F.\ Hegel \& Søn), 1903\\[10pt]
  \small Translated from the Danish}

% ---- original-page markers ----
% \opage{N} marks where printed page N begins, at the same joint as in
% transcription.tex (mid-sentence where the page turns mid-sentence).
\usepackage{marginnote}
\newcommand{\opage}[1]{\marginnote{\footnotesize\textit{[#1]}}}
\newcommand{\apage}[1]{\marginnote{\footnotesize\textit{[#1]}}}
\begin{document}
\maketitle
\thispagestyle{empty}
\tableofcontents

\mainmatter

% --- p. 85 ---
\piecehead{On the Good Will as a Power in Science}{\opage{85}\emph{R. Nielsen: Om den gode Vilje som Magt i}\\
\emph{Videnskaben}\\[2pt]{\normalsize[On the Good Will as a Power in Science]}\\[2pt](1867)}
In reply to Bishop Martensen's treatise \emph{Om Tro og Viden},
Professor \emph{Nielsen} has published, under the imprint of the
\emph{Gyldendal Bookshop}, \emph{Om den gode Vilje som Magt i Videnskaben}, a broad
and heavy-handed, a scornful and biting little book, written with an
overflowing zest for producing and full of good sallies. Its
merriment, however, has in places a character more Dutch than
Danish.

That the title is ironic, since the book sets out to show that in
science the good will counts for nothing, is a piece of information
that needs to be given to those who still remember that there was a
time when Professor Nielsen himself made the theologized conscience
the starting point of true science, and when it was proclaimed from
the Professor's own lips that, e.g., the freethinker's biblical
criticism, though \emph{impersonally true}, was without value, because
it had its ground in a \emph{personal untruth}. It is all the more
necessary to stress this, since fossil remains of this view still
occur to this very day in the productions turned out by Professor
Nielsen's present followers, so that these may see that they must
make some haste if they mean to keep up, and not, like their
predecessors, be left astern and lost along the way.

The book can be seen from two sides: it means to tear down and it
means to build up. It means to \emph{tear down} theology. What is
% --- p. 86 ---
\opage{86}theology? It is a discipline that stands in the same relation to
philosophy as alchemy to scientific chemistry and astrology to
astronomy. It has outlived its kindred spirits; it is one of the many
relics of the Middle Ages that are slowly crumbling away in modern
society. A fight against this dying thing has lost all interest; it is
a spectacle unworthy of the educated man, the working man, the
striving man, from which he would far rather turn his eyes away. How,
then, can Professor Nielsen go on regarding theology as a fit object
for an attack with sharp weapons? Who attacks a defenseless old woman
who can scarcely hold herself upright on her crutches? Who leaps with
drawn blade aboard a sinking hulk that is slowly going to the bottom
with every man and mouse aboard?

I know of no other answer than this: whom one loves, one chastens.
Professor Nielsen answers Bishop Martensen as he answered Pastor
Zeuthen; he has always had a little to spare --- at least a little
printer's ink to spare --- for the ordained, say what he will. % `havt en Smule ... tilovers for': both `had a soft spot for' and `had a little left over for'; the parenthesis plays on the second sense.
The only thing that might cause surprise is that Professor Nielsen, in
his zeal as an educator, has forgotten the fine old schoolmaster's
rule: If I were not angry, you should have a thrashing; for Professor
Nielsen is angry, and he lays on all the same, and so he lays on
wildly and now and then, out of sheer zeal, raps himself over the
knuckles.

A reader who has taken in the whole impression of the critical
research of the modern age cannot but regret that Professor Nielsen,
under the influence of a vehemence younger than his years, has geology
assign Paradise its place in the ``nursery'' and the botanist send its
trees ``to hell''; for such expressions will unpleasantly remind him
of the time when science was polemical because it was not yet able to
be critical, of the time when people mocked tradition and miracle
instead of understanding and explaining the psychological origin and
the historical development of miracle and tradition. That time, you
see, is past, and it cannot be Professor Nielsen's intention to call it
back. For that he has, on the one hand, far too thorough a cultivation,
and on the other, far too thorough a goodwill toward this very
theology from which he set out and toward which he still feels himself
irresistibly drawn. This is a point one must always hold fast to, or
one will feel bewildered at every moment. There are, e.g., passages in
the Professor's book by which one might otherwise
% --- p. 87 ---
\opage{87}let oneself be misled into the mistaken assumption that the
Professor's last word actually coincided with the freethinker's first;
so, e.g., the excellently witty piece that ends with the words: ``They
(the theologians) recognized that they could no longer go on talking
loudly about the supernatural facts of history, since the supernatural
facts were presumably not historical, and the historical ones not
supernatural''; but no one must let himself be bewildered by such
things. When Professor Nielsen here appeals to experience and to
``sound common sense'', which, as the serpent, gives theology to eat of
the tree of knowledge, he does not mean it seriously; his own
standpoint is not that of ``sound common sense'' but a quite different
one, which shall now be demonstrated as we pass on to what Professor
Nielsen means to \emph{build up}.

``Sound common sense'' sees, e.g., that there is exactly as much and
exactly as little reason to assume demons as to assume centaurs. A
scientific critic explains how belief in demons arose among the peoples
of antiquity. Professor Nielsen, finally, goes about it thus: he
% Danish prints a stray comma after `sig' (`bærer sig, saaledes ad');
% rendered as the sense, without it.
begins by ignoring the objection of sound common sense; next he grants
the critic that he is right ``from the standpoint of science''; but then
he himself adds to the latter's critique of knowledge a ``critique of
faith'', which obeys laws of its own and has nothing to do with the
critique of knowledge, is not crossed by it, does not collide with it
either in war or in peace, because any collision is impossible, since
the two critiques each obey an absolutely different principle. The
critique of knowledge has abolished demons, but at the same time
explained them; the critique of faith keeps them wherever their
existence ``makes visible the fundamental miracle of faith''; that the
whole question of demons is, for sound common sense, devoid of all
interest and settled in advance does not come into consideration. How,
then, do matters stand when the critique of faith comes to learn how
loftily the critique of knowledge, which keeps its papers in another
drawer of the brain, treats the critique of faith's own demons? The
secret is: it does not come to learn it. How so? ``For the moment the
critique of faith comes to learn it, it is no longer critique of
faith, but critique of knowledge.'' It is impossible that the river
Tajo can flow into Portugal, for the very moment it flows into
Portugal it is no longer called Tajo, but Tejo.

The critique of faith and the critique of knowledge, then, can as
little hinder or
% --- p. 88 ---
\opage{88}limit each other as faith and knowledge themselves, which
Professor Nielsen, calmly unconcerned about all objections, continues
to describe as ``absolutely heterogeneous principles''. The further
development of this is as follows, and most peculiar:

{\small
Consciousness is halved by relatively homogeneous factors, but not by
absolutely heterogeneous ones. When a consciousness is halved in thought, it
gets half a thought instead of a whole one: what is lacking is thought; if it
is halved in will, it gets half a will instead of a whole one: what is lacking
is will. On the other hand, one cannot in the stricter sense say that
consciousness is halved by thought and will, such that half of the thought
becomes will, half of the will thought [can one say this nonsense in any
sense whatever?], for thought can be transformed into will, and will into
thought, only in a highly figurative sense. [It is natural enough that the
figurative expression --- transformation --- must be taken in a figurative
and not in a literal sense]. If it is established, then, that faith can in no
way [a moment ago it was only: not in the literal sense] be transformed into
knowledge, or knowledge into faith, then it is also established that faith
can as little be halved by knowledge as knowledge by faith. If a
consciousness cannot be halved by friendship and mathematics, by love and
chemistry, still less can it be halved by ethics and astronomy, by religion
and geology.\par}

It is a drawback that precisely what was to be proved has here been
presupposed. That a true morality and a true religiosity cannot collide
with a true scientific spirit goes without saying. But the question is
precisely which morality and which religiosity is the true one, a
positive one or a free one? The question is, e.g., whether the ancient
religious representations of a small Semitic tribe, often enough at odds
with itself, from the most diverse periods --- representations that are
at every point interwoven with, and borne up by, its whole world view,
its ideas of the origin of the earth, its superstition, its countless
prejudices --- whether these do or do not contradict our geology and our
astronomy. They contradict them, in the same way as in countless cases
they contradict our morality and our religious feeling. But if Professor
Nielsen by no means succeeds in keeping the two hemispheres apart, he
succeeds far less still, in the end, in uniting them. For however sharp
the separation is claimed to be, a union must, after all, be
conceivable: ``The two absolutely heterogeneous principles meet, part,
and \emph{in a way} enter into negotiation with each other in the medium
of \emph{meaning}''. Meaning --- that is the saving category; on the % `Mening': meaning, sense, opinion; rendered `meaning' throughout the piece (also in the Genesis quotation, p. 89, whose `i den Mening' is Nielsen's illustration of it).
neutral ground of meaning the two absolutely
% --- p. 89 ---
\opage{89}heterogeneous ones negotiate with each other. The only party that
could have anything against such a negotiation would be the Danish
language; for it is indeed the philosophers' prerogative to teach
language some sense; but the two heterogeneous ones are related, after
all, as love and mathematics are; what is the name of the medium in
which love and mathematics conduct their negotiations?

Professor Nielsen elucidates by an example: ``What is told in the first
book of Moses about Paradise, about the tree of life, the tree of
knowledge, etc., is certainly told with the meaning that all of it
belongs to external reality, that the tree of life and the tree of
knowledge have their roots in the sensuous soil of the garden just as
surely as the vine or the fig tree. But in this meaning we cannot
possibly take it \emph{at the present stage of cultural consciousness}?''

Did we read that right? Cultural consciousness? Is Professor Nielsen
forgetting himself and his own assertions? What business has cultural
consciousness here, pray? Is the critique of faith to be grounded on
cultural consciousness, and faith nevertheless be absolutely
heterogeneous with it, not to be determined and not to be limited from
it? What a web of brain-spun fancy, what a muddle, out of which there
is no way out! Moral: do not give cultural consciousness your little
finger, or it will quickly take the whole hand.

``Those images,'' says Professor Nielsen, ``have not outer but inner
reality.'' That this is not new is known to everyone who knows anything
of the history of theology; that it is unintelligible is found by
everyone who tries to think it through. What is inner reality, if it
is to be different from poetic reality? --- ``The legend of Paradise,''
says Professor Nielsen, ``cannot possibly be myth, as surely as it
reveals to the eye of faith, in pure, clear strokes, one of the deepest
mysteries of the will.'' Does the myth of Prometheus perhaps not do the
same? And is this a distinguishing mark? When the legend is here denied
outward reality, it is a matter of indifference whether it is called
myth or not. That Paradise has a poetic reality which the ``Jura'' and
``graywacke'' do not have goes without saying and is denied by no one.

To the psychological union in the clime of \emph{meaning} there now
corresponds the essential (ontological) one in that of \emph{mystery}. The
development of this point is given in a language that regrettably recalls
the most unintelligible dialect of the Hegelian age.

{\small
% --- p. 90 ---
\opage{90}Faith and knowledge can be united only in an absolute mystery. The
absolute mystery is at once exclusive boundary and negative middle. As
exclusive boundary the mystery sets a separation between faith and knowledge
as between two regions. When the believer looks toward the center of faith,
and the knower toward the center of knowledge, what is true in faith becomes
untrue in knowledge, and vice versa. As negative middle, on the other hand,
the mystery unites the separated extremes and shows itself as the medium of
passage between the regions.
% Danish prints a comma after `Gennemgangsmedium' before capital `Idet';
% rendered as the sense, a full stop.
When the believer and the knower both at the same time direct their gaze
toward the mystery, the strife is settled and the extremes are reconciled.
Seen from the standpoint of knowledge, the miracle is now possible [yes, what
can a punctilious drill not achieve, where one looks to the right and to the
left at the same time!]. For while it is indeed impossible that a breach
should ever occur in the eternal laws of nature, which are identical with the
necessity of reason, yet with the immutability of the laws --- when the
intermediate determinations hidden in the mystery are understood --- the
manifestation of the power of will, of wisdom and of love, in which the
believer sees (!) the miracle, can very well be thought to subsist. Seen from
the standpoint of faith, the necessity of reason, with the immutable laws of
nature, can be united with the presupposition of the miracle. For while it is
indeed impossible that necessity should be able to hamper freedom, and the
laws of nature limit God's will, yet, when the intermediate determinations
hidden in the mystery are understood, it can indeed be thought that the laws
of nature, far from being broken, are on the contrary fulfilled in the very
highest sense every time a divine act of will is accomplished.\par}

This was good theology in the old days. Philosophical hints of
something of the sort are found, e.g., in the German clergyman
Bockshammer, who is cited by J. L. Heiberg in \emph{Prosaiske Skrifter} I, 110.
He concludes with a flowing torrent of words that contains the
self-same assertion. For more than an unproved and unprovable assertion
this is not, after all. Let us choose a miracle from the New Testament,
e.g., the scorching of the fig tree by the divine word of power that
punishes it for not bearing fruit. One need not, after all, be a thinker
to see that a miracle like this --- whose psychological origin,
incidentally, seems to be the misunderstanding of a parable as though it
were historical reality --- that such a miracle cannot be brought into
agreement with the laws of nature either in a plain and simple sense or
in any higher or ``very highest'' sense. That Most Highest Sense will
% `Allerhøjstsamme': the form in which Danish official style refers to the
% King (`His Most High Self'); the pun on `allerhøjeste' can only be hinted.
only waste its strength here.

The duel between Bishop Martensen and Professor Nielsen, these two
so practiced and accomplished champions, makes --- be it said with
sincere regret --- the same impression as a practice bout between two
equally poor fencers: neither of them can use his cutlass; neither of
them parries the opponent's blows; but
% --- p. 91 ---
\opage{91}for both of them the only point is who can deal the other the
quickest and heaviest blows. They strike like the old Norse berserkers,
with their shields on their backs, and one is given no occasion to
delight in great art or dexterity on either side. But if there is not
much art here, there is all the more politics.

It was good politics on the part of His Right Reverence, Bishop
Martensen, to leave Professor Nielsen's attack on his Dogmatics
unanswered; for dogmatics fare nowadays like other fortresses that a
hundred years ago seemed all but impregnable: they can no longer be
defended; their ramparts are going downhill, their moats are gutters,
and if the fortress will not surrender its keys, it does not matter,
since there are no gates.

But just as surely it was good politics on the part of the philosophical
champion to take the offensive and answer Bishop Martensen's polite
``Good day'' with a terrible ``Axe-handle''. Professor Nielsen's % `Goddag --- Økseskaft': the Danish byword for an answer that has nothing to do with what was said (`Goddag Mand! --- Økseskaft!').
main guard is no good for defense. Professor Nielsen does have one man
on guard, but he weeps --- with the one eye and laughs with the other;
his right hand knows not what his left is doing. He advances on two
heterogeneous legs, a leg of flesh and a leg of wood, on two absolutely
split legs, cleft like one of Uffe's Saxons. His head is a Janus head: % `Uffes Saksere': Saxo's Uffe the Meek, who fought two Saxon champions alone at the Eider.  `Kødben' (leg of flesh) is also `meat bone'.
with the one mouth he proclaims the higher fulfillment of the law of
nature through the miracle; with his other throat he gives the honor
to cultural consciousness and sings a well-meant hymn of praise to
David Strauss. This last is perhaps the oddest among many odd features
of this feud. The congregation is admonished by Professor Nielsen, far
from taking offense at Strauss's audacity, to appreciate the help that
this man, by his ``consistent and resolute'' procedure, has directly
rendered faith, in that he has founded the critique of knowledge to
which Professor Nielsen, by his projected critique of faith, means to
supply the necessary counterpart. Bishop Martensen must take a scolding
for his orthodox condemnation of Strauss. ``If Strauss is to be
excommunicated, then there are many who must be excommunicated'' [yes,
% The Danish never closes the quotation opened at `‹Skal Strauss';
% the English closes it before the bracket.
why not, that is after all a good Christian doctrine!]; ``if Strauss is
an Antichrist, then there are many Antichrists in the camp'' [yes, what
orthodox man would ever doubt it!] The Janus head turns to both sides;
it stands like Don Juan between the peasant girls and has fair words for
% --- p. 92 ---
\opage{92}both parties. The Grundtvigians get their share of them. \emph{Dansk
Kirketidende} receives public praise; Mr. N. Lindberg is commended ---
for what, would you think? --- for \emph{logical-categorical certainty}.

The Fool says in \emph{Aladdin}: ``I am not fooling, it stands, upon my faith,
% `Tro' is set in the Danish with a wrong-fount o; no bearing on the English.
in the letter.''

If anyone will take the trouble to look up the article in
question, he will, it is said, be able to find in its notes interesting
samples of the certainty that Professor Nielsen singles out for praise,
and so be able to convince himself with his own eyes, not of the way in
which a church paper takes part in a scientific polemic, for who looks
for science in church papers, but of what, in Denmark, a philosopher
can praise. Professor Nielsen does not forget his Pappenheimers; his
Pappenheimers will not forget him either.
% `Pappenheimere': Schiller, Wallensteins Tod III.15, ``Daran erkenn' ich
% meine Pappenheimer'' -- one's own people, whom one knows.
\piecerule

% --- p. 93 ---
\piecehead{Nielsen and Grundtvig}{\opage{93}\emph{Gb. (Rudolf Schmidt):}\\
\emph{R. Nielsens Filosofi og den Grundtvigske Anskuelse.}\\[2pt]{\normalsize[R. Nielsen's Philosophy and the Grundtvigian View]}\\[2pt](1867).}
Under the signature \emph{Gb.} an author who stands equally close to
Grundtvig and to Nielsen has published a book of 124 pages, entitled
\emph{R. Nielsens Filosofi og den Grundtvigske Anskuelse}, a book whose
form bears witness to much stylistic maturity and cultivation. Its
standpoint is this: to side with Nielsen unconditionally, both where he
attacks and where he is the object of attack; to assert the genuineness
of Grundtvig's genius and the importance of his person and his doctrine,
while at the same time critically setting apart some of the great deal
that is loose and repellent and bound up with Grundtvig's whole
activity; and finally, at the end, as well as it can be done, to make
Nielsen and Grundtvig rhyme. The worth of the book's contents differs
greatly from one part to another. It can be stated most briefly and
pithily thus: the larger part, the section on Grundtvig, which is
written with a respectable striving for impartiality and rests on an
uncommonly wide reading in Danish literature, contains much that is true
and to the point; but the two short sections on Nielsen, which begin and
end the book, are downright wretched. They are not merely empty of
thought but written in an impertinent and indecent manner. For it is
impertinence to meddle in a serious dispute in order to deliver, with
lofty condescension, one's humble opinion as though it were a verdict of
the Supreme Court, when one is unable to furnish a single contribution to
the dispute or to refute a single one of the arguments
% --- p. 94 ---
\opage{94}one brushes aside, and contents oneself instead with the empty
phrase that the attackers have shown themselves to be ``without the power
of appropriation in the face of the new.''

Everything that has so far been brought forward against the new doctrine
is disposed of with the talk that, if one acknowledges Nielsen's uncommon
gifts, one cannot possibly credit him with such errors and contradictions
as are laid at his door. A greater perversity than the one underlying
this answer, repeated often enough, cannot be conceived. Suppose a man is
to go through a difficult set of accounts. This examination is regarded
as a formality! for those around him hold it as an article of faith that
these accounts must be correct; so they were always found by the earlier
auditors, and if any of these raised difficulties, the matter was hushed
up. The man would come into conflict with his own trust in the
authorities and fall out with all respectable people if he did not
arrive at the given total; and so, happily, he does make it tally.
Unfortunately, someone now discovers that this happy outcome is owed to
the circumstance that at one place in the accounts 2 and 2 have by
mistake been added up to 5. He points this out, and receives, one after
the other, four answers: \emph{a}) This objection rests on a
misunderstanding. \emph{b}) You must surely be able to tell yourself that
it is impossible that a master of figures like this one \emph{could}
make such a mistake. \emph{c}) What you say is abominable and
outrageous; it undermines all confidence in the established order and
the whole childlike trust in the authorities without which everything
falls apart. \emph{d}) What a threadbare and hackneyed objection! Truly a
most novel and remarkable discovery, that 2 and 2 make 4!

\emph{a}) The objections rest on misunderstanding. The answer is simple.
How long do you think people will put up with this evasion? If we
misunderstand, then open your mouths and speak plainly, so that one can
understand you, instead of talking as if you had porridge in your
mouths! It has come to this: it has now become the fashion in
philosophy, as it has long been in politics, that the highest art of
every pronouncement is to conceal where one is really heading, to
express oneself in riddles, to say yes and no like a diplomatic dispatch
or a speech by Napoleon III.

\emph{b}) Can it be that Professor Nielsen does not know that 2 and 2 make 4?

That does not concern us and is a matter of complete indifference to us;
but we point to the fact that here, on this particular point, he
% --- p. 95 ---
\opage{95}has let five pass for even. % `har ladet 5 være lige': the idiom `lade fem være lige' (let five be even) means to let a thing slide; Brandes takes it literally, of the sum. The pun is kept as far as the English allows.
Can anything more ridiculous be conceived than that, in the face of
the specific demonstration, people go on appealing to the man's
abilities? Contradictions have been demonstrated point by point in
Professor Nielsen's doctrine. There is no weapon against logic but
logic. But we ask logically, and you answer us lyrically. It is as if
the famous Bog Wife had become the Valkyrie of the whole camp. % `Mosekonen': in Danish folklore the mist rising over the marshes is the bog wife brewing; i.e. the camp's war-maiden is a maker of fog.
People come forward with presentations that go by leaps, with fanciful
word-paintings, parables, shadow pictures, everything except proof,
refutation, logic; for of these there is not even a shadow here.

\emph{c}) The objections are out to take the life of the ideals (see the
review of this book in \emph{Fædrelandet}); with dry common sense they
would rob faith of its ground and enthusiasm of its source. We answer:
% `til Hekkenfeldt': Hekkenfeldt (Mount Hekla) is the Danish folk
%   euphemism for hell, `to blazes'; kept, since the opponent quotes the
%   phrase back at Brandes (p. 98).
A fig for such ideals, and to Hekkenfeldt with an enthusiasm that creeps
into hiding in the bolt-hole of \emph{existence} when science challenges it
to come out into the sunlight of the spirit. Keep to the facts and refute
our arguments! Let those big words, the ideals, enthusiasm, which are
posted up like placards and impress the mob, go, and set proof against
proof. Do not boast, then, of your exact method, which is supposed to
make your chain of inferences as reliable as the mathematical one, when
everything you say has a merely belletristic and in no way whatever a
scientific character. The day of belletristic truths is over. Woe to him
who cannot bear the naked truth; woe to him who sees the task of
philosophy in pasting the fig leaf on it!

\emph{d}) What is brought forward against Professor Nielsen is old. Yes,
it is old that a thinker must not contradict himself; yes, the truth is
old, old as everything eternal, but what a complaint to come out with! If
there is anything paltry in the youth of Denmark, it is its inclination
to inferences like these: ``This is new, therefore true; this is new,
therefore ours; this is new, therefore that around which we must
throng.'' The novelty-mongers do not hear the call of science: ``You
% PRINTER'S ERROR: the Danish closes the quotation after `stimle' with a
%   turned (opening) guillemet; the English closes it.
fools! at least test this new thing first!'' ``No,'' comes the answer,
``let us hurry, let us hasten to get deep into the crush and push aside
anyone who goes against the current!'' And does one really think that
anyone will be deterred by that? Does one really think that anyone who,
in questions like these, so important, so far-reaching, so agonizing,
alone with himself but face to face with the spirit of science, has
wrested from it a conviction as living and as powerful as the spirit
itself, cares much about these currents of the age on the
% --- p. 96 ---
\opage{96}surface, or about those who feel safe and strong only when, in
the midst of the crush where no one can move arm or leg, they are borne
along without a will of their own in the wake of some admired figure? We
have no claque in our theaters; are we now to get a claque in our
literature?

No, it is impossible to give up the hope that the young people of
Denmark, for all their indolence, will try to arrive at an independent
judgment. Let them feel their way forward; let them resolve, for once,
to want to see with their own eyes! Fortunately we are not, after all,
intellectually bound to the soil; there are ways out for anyone who feels
% `aandeligt stavnsbundne': bound to the soil like the Danish peasant under
%   the adscription (Stavnsbaand), abolished in 1788.
that in one branch of our literature we lag behind the rest of Europe.
Let the young study foreign thinkers, and it would truly be strange if
they did not discover that these declamations are not refutations; let
them study the more recent theism, e.g., and they will see the connection
between Professor Nielsen's endeavors and the European reaction against
what is taken to be a one-sided pantheism; let them only acquire a
horizon, a wide horizon, and they will discover how far from unique
Professor Nielsen's ``immortal'' treatment of the problem of subjectivity
is, the treatment that the pamphlet sees in connection with Old Norse
mythology.

What is at work here in leading many astray is the following reasoning:
When everyone must concede that the religious problems have been treated
in our own literature (first and last by Kierkegaard) with a range and a
depth hardly matched, perhaps, in any European world of books, how can
anyone believe that in any other we might find greater clarity about
these questions than here at home! This reasoning is very seductive, but
for all its apparent force it is quite hollow. For it is by no means
always the case that the most violent crisis brings about the best and
truest result. No one can deny that, set beside the French Revolution of
1789 and the French revolutions that followed it, those gigantic
exertions to win freedom by fighting for it, our revolution of 1848 ---
the long procession of respectable men who strolled arm in arm to
Christiansborg --- though not lacking dignity in itself, is pitiful,
meek, and of no account. And yet --- where is there more freedom now, in
France or here? But in the religio-philosophical sphere the relation is
% `omvendt forholder det sig': the case is reversed in that the violent
%   crises are now Denmark's and the lesser ones France's; the moral is
%   the same as for politics.
reversed. How slight France's religious crises in this century are
beside ours! And yet --- where is there more
% --- p. 97 ---
\opage{97}truth, in French religio-philosophical criticism or in our own
productions in a similar direction? Anyone who has studied French
literature even to some extent will hardly doubt to whom in this field he
must hand the palm. For Kierkegaard is indeed --- and there can be no
question here of anyone but him --- an immense spirit, one of those who
are born scarcely once in a century; but Kierkegaard is the Tycho Brahe
of our philosophy. He is great as Tycho Brahe was great, but like the
astronomer he does not dare, in his blind reverence for authority, to
place the center of our system in its sun. In astronomy this sun is a
body; in philosophy it is called reason. It seems a tragic fate attached
to Denmark that its greatest discoverers should be mistaken about what is
central. But just as posterity was not prevented by Tycho Brahe's errors
from turning his work to account, just so little does Kierkegaard fall
along with the positive element in morality and religion. That element
can be separated from his presentation, and infinite treasures will still
remain. Those for whom this is a closed book, as it is for Professor R.
Nielsen or for Mr.~\emph{Gb.}, are fighting only against progress, and
striving, though in vain, to keep the true new from breaking its way
through. But this new is irresistible; for it bears reason and freedom on
its shield.
\piecerule

% --- p. 98 ---
\piecehead{A Reply}{\opage{98}\emph{A Reply}\\[2pt](1867)}
The author of the pamphlet \emph{R. Nielsens Filosofi og den Grundtvigske
Anskuelse} has placed in yesterday's \emph{Berlingske Tidende} a rejoinder
to my attack on his book, which, without bringing forward anything
whatever against the four points I raised, ends with the words that
``those who really believe that there is a cause at stake here will not
fail to see that he (i.e.\ I), stung by the goad of self-importance,
merely dances a St.\ Vitus's dance around the great tasks.'' Such things
it is enough to quote; there is no need to answer them.

As regards the accusations and questions that my honored opponent raises
besides, it seems to me, however, that I owe an answer, not to him, but
to the readers of \emph{Dagbladet}. He reproaches me with having pointed
to the French religio-philosophical criticism, ``which newspaper readers
do not know.'' I do not see what harm there would be in that, even if it
were so; but of course it is not so. In our day newspaper readers, as is
well known, do not form a caste of their own; practically all educated
people read newspapers. My honored opponent is himself an example. No
doubt Mr.~Gb. is a most cultivated man (see, e.g., the sample of his
% The Danish prints `Hr. Gb en' without the stop after `Gb'; the English
%   sets `Mr. Gb.' as elsewhere in the piece.
style above), but why should \emph{Dagbladet} not be so fortunate as to
count among its readers several men just as cultivated, men who among
other authors know Renan, whom Mr.~Gb., as we know, knows so well?

My opponent next accuses me of having consigned to Hekkenfeldt religious
ideas ``that have borne mankind up through eighteen centuries.'' That
would undeniably be an ugly trait. But is it really quite certain that I
% --- p. 99 ---
\opage{99}have acted so? Is it certain that the ideas I would have
banished have for eighteen centuries \emph{borne human beings up?} (and
that, surely, is far less essential than whether they are fit to bear
mankind up today). The feuilleton that happened the other day to
accompany my article, the fragment on the witch trials of the Middle
% `Føljetonen': the literary section printed below the line in the
%   newspaper (Dagbladet), which that day carried the piece on witch trials.
Ages, which were engendered precisely by the supernatural ideas that I
dispute, makes a good contribution to answering this question. In that
age, when the supernatural really was alive in human consciousness, the
views that nowadays people are content to expose to hatred were punished
with the heretic's stake. --- A few words I could say in defense of the
likeness between religious and political crises I shall spare, since my
opponent has so puffed himself up that he \emph{forbids} me \emph{to speak}
of it.

There remains, then, only one point, the question repeated again and
again, of whose crushing force people are so convinced: How can I
acknowledge Kierkegaard's greatness and yet be convinced that he was
mistaken about what is central? I could answer this question as I have
answered it before, with the words that only very dependent people
understand by admiring acknowledgment nothing other than simple adherence
through thick and thin; but since the question is brought up anew, I add
the following brief account of how I have come to this result: I
discovered fairly early that Kierkegaard was not infallible. I
discovered, e.g., and I suppose many have had the same experience, that
Kierkegaard was an adherent of political absolutism in Denmark; that in
the political sphere, with a remarkable narrowness in acknowledging the
inherited merely because it is inherited, he misunderstood his time and
fought against development and progress. I saw that his political
prophecies, of which he himself is so proud, did not come true at all. I
was thus led to observe that he was an absolutist in the religious sphere
as well: that to break with tradition, to desert one's parents', to abandon
a father's way of seeing was to him not merely a horror such as is surely
felt, though in differing degrees, by everyone who is compelled
% PRINTER'S ERROR: the Danish prints `fil at se' for `til at se'; rendered
%   as the sense.
to look filial piety in the face, and thus not merely a dread that it is
man's task to overcome, but a \emph{spiritual} horror of the same kind as
% PRINTER'S ERROR: the Danish prints `at samme Art' for `af samme Art';
%   rendered as the sense.
the one that nature implants to keep man from committing a crime against
nature. I believed, and
% --- p. 100 ---
\opage{100}believe, that he is wrong in this; I hold nothing sacred merely
because it has been handed down; for being handed down makes a thing
neither rational nor true. I regard it as beyond doubt that on this point
one can learn many times more from the religio-philosophical writers of
Germany and France than from him.

It is well known that Kierkegaard's works are distinguished by a rare
consistency and unity. But it is, if perhaps not generally acknowledged,
hardly the less certain for that, that the attempt he made at an advanced
stage of his career as an author to demonstrate the inner coherence of
all his writings is, on the one hand, very forced, and on the other
cannot rule out that these works keep their value in and for themselves,
without regard to their place in the system of the writings. Is
Kierkegaard on this point anything other than what he would least of all
wish to be, and what he hated most of all --- a systematizer who laces the
free fullness of life into the stays of a preconceived, or rather in part
an afterwards-conceived, fundamental thought? He himself is infected by
the systematic mania that he combated in his age. But even if it were
otherwise, even if everything Kierkegaard has written really did aim,
from first to last, at presenting the specifically Christian in the form
in which he conceived it, would anyone of another mind therefore be
forbidden to take delight in, and make his own, the numerous productions
of his that have nothing to do with the paradoxical religiosity? Seen
from an alien standpoint, is the essay on Don Juan any less inspired,
\emph{In vino veritas} any less brilliant, the essays on marriage any less
captivating, or the presentation of non-Christian religiosity (in
\emph{Afsluttende Efterskrift}) any less significant and less instructive?
What an absurdity!

But there is no need at all to stop here. In a great many cases where
Kierkegaard, in his own opinion, presents the paradoxical (as in
\emph{Kærlighedens Gerninger}), this paradoxical element is not contrary to
reason at all and can be translated into the language of the human; in
other cases, e.g.\ in \emph{Sygdommen til Døden}, the positive comes in
toward the end without being derived. It is impossible, starting from the
concept of despair as Kierkegaard at the outset lays it down in purely
human terms, to bring despair, and with it sin, to the high point in
rationalism's view of the Christian. There is, then, much in this that a % Doubtful construal: `at faa Fortvivlelsen og dermed Synden til at naa Højdepunktet i Rationalismens Betragtning af det Kristelige'.  Rendered literally, keeping the Danish ambiguity: `i Rationalismens Betragtning' may qualify `Højdepunktet' (the height sin reaches in rationalism's view of Christianity) or give Brandes's own standpoint (`in rationalism's view, it is impossible ...'); the first translator read the former, the reviewer the latter.
% --- p. 101 ---
\opage{101}worshipper of reason can make his own, and much else that he will
cut away as an excrescence.

Such a separation of what one regards as the lasting in Kierkegaard from
what one considers the transient is something that simply cannot be
avoided at all. If one will have no such separation, then let one take
the consequences, and with Kierkegaard in his last public stand reject the
whole life of the church and the whole life of human society, marriage,
etc., which is surely the last thing the half-Grundtvigians would want.
% `Halv-Grundtvigianerne': a jab at Gb. and those who, like him, stand
%   half with Grundtvig and half with Nielsen.
Let no one insist that he is at one with Kierkegaard on what is central
when he himself will at any price break off the finest and sharpest point
in which his whole endeavor culminates. All or nothing, that cry so
popular in our day, is a watchword that, where acknowledgment is in
question, signifies nothing other than an attempt on the life of all
criticism.
\piecerule

% --- p. 102 ---
\piecehead{Mr. 99 of Fædrelandet}{\opage{102}\emph{Mr. 99 of ``Fædrelandet'' (Bjørnstjerne Bjørnson)}\\[2pt](1867)}
The honorable gentleman who has come out against me in
\emph{Fædrelandet} under the mark 99 possesses, among other engaging
qualities, an uncommon effrontery. It sounds incredible, but it is
true, that Mr. 99 has not even read for himself the book over
which he has taken up arms. The quotations he prints from this
book are nowhere to be found in the work; what he puts forward
as the book's words are \emph{sentences of mine}. Likewise:
what he quotes from \emph{Fædrelandet} is not the paper's words
but \emph{mine}. He quotes me alternately in small type (as proof
of how presumptuous I am) and in large type (in the belief that
he is setting my opponents' words against mine); it seems to me
that this little fact is in and of itself enough to show what
kind of opponent he is.

Like his kindred spirit in \emph{Berlingske Tidende}, 99 wisely
skips over all the main points and throws himself upon the
quite inessential comparison I made. I said: it is not always
the most violent crisis that brings about the best and truest
result; thus, e.g., in the middle of this century we Danes in
the political sphere, and the French in the religio-philosophical,
have taken a great step forward without violent crises, and
conversely the political revolutions in France and the
theologico-philosophical movements in Denmark have not led
directly to freedom. This was the comparison, and it seems to me
% --- p. 103 ---
\opage{103}that it is clear. One can of course ride it to death, as one
% `hidse Aanden af den': to overdrive it, as a horse is driven till it
% drops; construed so.
can any comparison; but one has no right to want it stretched
further than I have stretched it; for otherwise it is truly easy
enough to make it limp: ``In the French Revolution blood was
shed, in our theological disputes none was shed, therefore
etc.''; ``Napoleon III put an end to the last revolution, our
science has no Napoleon III, therefore etc.'' That our political
movement presupposes the French one, while the French know
nothing of Kierkegaard --- German literature they know very well
--- is quite irrelevant to the comparison; but for the rest, this
conception of the reason for the relation I mentioned is as
shallow as it could possibly be. The reason is that the Romance
and the Gotho-Germanic stocks have quite different capacities
and tasks; the Romance stock, represented by the French,
conceives and develops morality socially, the Gotho-Germanic
stock conceives and develops it personally. The former always
makes the beginning socially; from the latter the reforms
proceed in the personal sphere; and each appropriates the
results of the other's efforts. In the direction in which each
of the stocks is least developed, it receives the results from
the other without much struggle, whereas the more developed one
may fight itself so weary that for a time it lies exhausted and
is unable to draw the consequences it has itself prepared. Mr. 99
stamps this view as ``nonsense'' and the comparison as an
``unseemly blunder''. There can
be so great a gulf between the development of two people that
the words of the one, however clear and well-founded they may
be, must necessarily look meaningless to the other.

If any proof were needed of my remark that here people will
not refute but declaim, my opponents have now supplied it; for
with what have they answered me, if not with insults, with
high-handed pronouncements and with appeals to national feeling?
The one \emph{forbids} me to call Kierkegaard the Tycho Brahe of our
philosophy --- something, however, in which I intend to persist,
with or without his permission. The other would have me denied
all use of pen and ink whatever, at least in the daily papers. O
noble man! Yes indeed, it would be splendid if one could get rid
of the restless spirits in this way and sit so snug in one's own
little soap-cellar! Certainly it is easier to gag than to
% `Sæbekælder' (soap-cellar): a basement chandler's shop, the image of
% snug, petty, shut-in comfort; kept literal.
% --- p. 104 ---
\opage{104}refute; but alas, those golden days of the Middle Ages, when such
things went so smoothly and so easily --- they are over, those
feast-days\fnmark{*)}{An
interesting parallel, as regards the kinship between the
astronomical relocation of the center and the philosophical
relocation of it to revelation, is afforded by the following
quotation: ``\emph{Pascal} somewhere corrected a sentence of his
own, took it back and recast it, in order to make the sun turn
around the earth and not the earth around the sun. This great
spirit, which on this point still harbored a remnant of
superstition, shrank back from the truth of Copernicus.''\endgraf
% Brandes quotes Sainte-Beuve in Danish; translated from his Danish.
% The source line is kept verbatim, as printed.
Sainte-Beuve: Causeries du lundi III S. 254.}.

But I come to the most wounding attack that has been directed
at me. It is this: that ``the man who can feel and speak'' as I
do ``has no view in common with other Danes, has no Danish
consciousness''. My view is --- not illogical, but ---
\emph{un-Danish}. What is it, then, that I have said? That
Kierkegaard, in my conviction, went wrong in what is central,
i.e. in what he took to be central, the doctrine of the
paradox --- not, of course, in what I take to be central in his
writings, their purely human content. Is it a crime against one's
fatherland to say this? Has this pronounced ``the death sentence
on Danish culture''? Have we perhaps acquired a
% Danish prints `Kultur›?›', with a second closing guillemet after
% the question mark; the English closes the quotation once.
philosophical authority into the bargain, alongside the religious
ones? Kierkegaard must be right. Why? Because he is Danish. And
Martensen, his opponent? Right, because he is Danish. And
Grundtvig, his antipode? Right, because he is Danish. And Tycho
Brahe against Copernicus? Right, because he is Danish. Yes, it is
well enough known that this last Grundtvigian inference has
already been drawn, with sublime contempt for all sound reason.
Down with him who dares to say that ``Danish culture'' is not in
% Danish prints `af den ‹danske Kultur›' for `at den'; rendered as the
% sense (`that').
every respect the best, soundest, truest on earth!
% `Nittengryn og Nioghalvfemser': the pun on 99 (nioghalvfems) and
% nitten (nineteen) cannot be carried; Nittengryn kept as a name.
% `stavnsbunden': bound to the soil, as under the old adscription
% (Stavnsbånd) of Danish peasants.

``Will \emph{Dagbladet} henceforth keep contributors,'' exclaims
Mr. 99 --- the expression ``keep'' is itself telling ---
contributors ``who repeatedly encourage the youth of Denmark to
come and see with their own eyes, if they do not wish to be
called intellectually bound to the soil''? Yes, is it not dreadful
to urge the Danish youth to see with their own eyes? No, let us
stay in the rut we are in now, let us all together become
Nittengryns and Ninety-Niners; for what does it not say in the
Fox's \emph{Morals} in Holberg, Art. 17 and Art. 18: ``\emph{Scruple}
not in religion;
% --- p. 105 ---
\opage{105}but hold thee blindly to the reigning faith so long as it is on
the throne, for it is most seemly and safest to follow the
fashion, in this as in dress. --- Trouble thee not to search out
the truth; but believe what thou art bidden to believe; for the
first is both dangerous and troublesome, the second, on the
other hand, safe and comfortable; yea, it often paves the way to
\emph{canonization}''.
\piecerule

% --- p. 106 ---
\piecehead{R. Nielsen's Doctrine of Faith and Knowledge}{\opage{106}\emph{Heegaard: Professor R. Nielsens Lære om}\\
\emph{Tro og Viden}\\[2pt]{\normalsize[Professor R. Nielsen's Doctrine of Faith and Knowledge]}\\[2pt](1867)}
There has just appeared \emph{Professor R. Nielsens Lære om Tro og
Viden}, a
critical treatise by Dr. \emph{P. S. V. Heegaard}. It makes no
secret of what it bears on its shield; for on the very cover it
carries the motto: ``A philosophy that is in open contradiction
with sound understanding is false.'' The book is a literary
curiosity, not so much for its content --- for part of what it
presents has been said before --- as for the surprise it
occasions, to see \emph{this} view expressed from \emph{this}
quarter; a pleasant surprise, to be sure, but nonetheless, as
Nonpareil says in \emph{Recensenten og Dyret}: la surprise
surprenante des surpris.

For the author, as is well known, was one of Professor
Nielsen's most faithful; he has been the Professor's pupil,
adherent and secretary, and now --- alas, how changed! Where now
is the time when Dr. Heegaard, as the highly trusted disciple,
lay on the Master's breast and rose only to chastise with severity
the \emph{rash} attempt of an \emph{immature} youth to tear down a
certain solidly founded doctrine? Yes, where are the snows of
% `hvor er den Sne, som faldt ifjor?': the Danish form of Villon's
% refrain; given in its familiar English form.
yesteryear? And yet Dr. Heegaard was right then. Was it not
perhaps \emph{precipitate?} Why should the young people hurry so,
why come two whole years too early? Could they not easily have
waited a little and taken Dr. Heegaard along? After all, he only
wanted a little trip abroad first (see the last page of the
book). And was it not
% --- p. 107 ---
\opage{107}an \emph{immature} undertaking to sound the attack before the
crack troops had arrived, to charge without waiting for so brave
and valuable a reinforcement? Yes indeed: for how did it go?
Alas, miserably; the ranks wavered, the attackers were done for;
then, at the decisive moment, Dr. Heegaard arrives like Desaix at
Marengo. Straight out of Professor Nielsen's own camp he breaks
forth all at once, with flying colors and full band, leaving
behind nothing but a trumpet, which, it is said, is now vacant.
He hastens over to the opposing side, he falls upon the
Professor's flank; it is in deadly earnest, and one comes over
quite solemn at such a spectacle. Happily, a little jest and fun
mingles with the solemnity of the procession and the horrors of
war.

For example: what has offended him and provoked him to deliver
his bloody critique, says the author, is the claim that Professor
Nielsen, as a mind of the first rank, has brought about a great
philosophical breakthrough. This is of course a jest; for
as is well known it is the author himself who earliest and most
eagerly set this claim down in writing. Further: it is most
unfortunate, Dr. Heegaard thinks, that Professor Nielsen, who
began as a ``pure Hegelian'', later became Hegel's decided
opponent. Another jest --- namely, a jesting illustration of the
proverb: One should not speak of rope etc. Finally: consider
\emph{Grundideernes Logik}, says the author; the best thing would
be for Professor Nielsen, who is after all unable to complete the
building, to leave it standing unfinished at once, as an
``interesting ruin'', and so as a kind of Marble Church. ``Marble
Church yourself,'' answers Professor Nielsen, calmly pointing to
the first part of the far-famed \emph{Indledning til den
rationelle Etik}. Now is this not a merry situation?
% `Man skal ikke tale om Strikken osv.': the proverb runs on `in the
% hanged man's house'. The Marble Church (Frederik's Church,
% Copenhagen) stood as an unfinished ruin from 1770 until the 1870s;
% Heegaard's own Indledning til den rationelle Etik had reached only
% its first part.

But joking aside, it is clear enough that precisely the mixed,
half pathetic, half comic position in which Dr. Heegaard finds
himself as the author of this book greatly increases the weight
of his critical pronouncement. With what irresistible power must
the truth not have shown itself to him, since
% Danish prints `ham. siden' (a full stop before lower-case `siden');
% rendered as the sense, one sentence.
he has made so great a sacrifice to it! His critique accordingly
bears witness first and foremost to a truth-loving mind, resolved
no longer to let itself be dazzled, not to be content with
flourishes or riddles, to lay bare the web of subtleties and
fallacies on which Professor Nielsen's so-called philosophy of
religion rests. The critique is as able as it is serious, and as
sharp as it is dogged
% --- p. 108 ---
\opage{108}and merciless. It tears apart the flimsy cobweb of the brain in
which so many weak souls have been caught like flies, and it is,
% `Hjernespind' (a figment spun by the brain): `cobweb of the brain'
% keeps the image of the flies.
as a whole, so decisive and definite that it would now, after all, seem bound to
become clear even to the slowest head that Professor Nielsen has
lost the battle and that his doctrine has gone the way of all
flesh.

What is new in Dr. Heegaard's treatment of the question is chiefly
this: that instead of showing how Professor Nielsen's doctrine
conflicts with the truth in and for itself, he proves it by a
detour, through a demonstration of how Nielsen, in every branch
of his science, comes into contradiction \emph{with himself}
through his reason-defying conception of the relation between
faith and knowledge. For a critique of this nature Dr. Heegaard,
with his special qualifications, with his exhaustive knowledge of
Professor Nielsen's entire system, was admirably equipped, and it
leaves very little to be desired. Here and there, to be sure,
there is some unclarity (as, e.g., in the conception of the
Mystery as a category), and now and then a somewhat perverse
interpretation of the opponent's words (as where the Mystery is
equated with the Meaning); but these small defects vanish beside
the thoroughness and sharpness that otherwise distinguish the
argument line by line.

The book is \emph{purely} critical; it barely so much as hints at
a standpoint of its own; the author seems to have only just come
out of the Nielsenian school without yet having gained a
foothold himself. It cannot be denied that this circumstance
gives the book and its style something dry and coldly cerebral;
there is in the author's personality no primitiveness, none of
that native originality which, with the sureness of instinct or
of genius, seizes what presents itself to it as truth; as a
writer he seems to lack the best thing of all: the productive
gift which, after the habit of barren natures, he would deny to
\emph{the age}.

That is why it surprises us that he does not value this capacity
more highly in his former teacher. For much can be said against
Professor Nielsen, and much ought to be said against him: he has
entangled himself in a religio-philosophical maze, and as a
thinker he is no systematician at all; but he has the luxuriantly
fertile mind, the eye of genius, the rare originality, which
are just as essential capacities in a thinker as the strictly
scientific cultivation of thought and a sharp and reliable
understanding.
% --- p. 109 ---
\opage{109}One misses in Dr. Heegaard's book a proper acknowledgment of
this. The gratitude which undoubtedly dwells in the author's mind
toward his so highly gifted teacher could have found fuller
expression; it ought not to have been squeezed into a brief
declaration, but ought to have pervaded the whole development and
made impossible the irritated tone in which his critique speaks
more than once.
\piecerule

\end{document}
